\documentclass[12pt]{article}
\usepackage[utf8]{inputenc}
\usepackage[T1]{fontenc}
\usepackage{hyperref}
\title{Simbo Language Resources for K-12 Curriculum Development}
\author{Compiled by Claude Sonnet 5.5 (claude-sonnet-5-5) for Andr\'as Kornai}
\date{}
\begin{document}
\maketitle

\section{Introduction}
Simbo (ISO 639-3 \texttt{sbb}; Glottocode simb1256) is an Austronesian language of Simbo Island in the Western Province of the Solomon Islands. Its speakers call it Mandegusu (also Madeggusu). Ethnologue gave 2,701 speakers in 1999. It is not known to be taught in schools. There is a Simbo dictionary and some older descriptive work, but no full reference grammar.

\section{Primary Educational Resources}
I found no Simbo primers, textbooks or school programmes. The nearest relevant item is a University of Melbourne project on the Kulu Language Institute on Ranoqa Island (Western Province), which teaches reading and analysis of local languages. It does not concern Simbo specifically: \url{https://arts.unimelb.edu.au/school-of-social-and-political-sciences/our-research/research-projects/learning-through-language-in-solomon-islands}

\section{Linguistic and Reference Materials}
\begin{itemize}
\item Simbo overview (Surrey Morphology Group, Northwest Solomonic): \url{https://www.smg.surrey.ac.uk/languages/northwest-solomonic/simbo/overview/}
\item Simbo dictionary by Bill Palmer, Version 3.0 (2007), compiled from secondary sources and Christine Dureau's 1990s fieldwork: \url{https://www.smg.surrey.ac.uk/languages/northwest-solomonic/simbo/simbo-dictionary/}
\item Earlier Version 2.0 of the dictionary, drawing on Hocart's 1908 texts and Tryon and Hackman's wordlist: \url{https://www.smg.surrey.ac.uk/languages/northwest-solomonic/simbo/simbo-dictionary/simbo-dictionary-version-20/}
\item Surrey bibliography for New Georgia languages: \url{https://www.smg.surrey.ac.uk/languages/northwest-solomonic/bibliography/new-georgia/}
\item Lanyon-Orgill (1969), \emph{The Language of Eddystone Island} (Glottolog record): \url{https://glottolog.org/resource/reference/id/112542}
\item Handwritten Simbo-English vocabulary compiled by Rivers, Wheeler and Hocart, held in the Royal Anthropological Institute archive (MS 217): \url{https://www.therai.org.uk/archives-and-manuscripts/manuscript-contents/0217-rivers-william-halse-rivers-and-others-ms-217}
\item Verbix page for Simbo: \url{https://docs.verbix.com/Languages/Simbo}
\end{itemize}

\section{Religious and Community Resources}
I found no Simbo scripture, hymnal or community organisation in search results. Roviana was historically a church lingua franca in the Western Province but is now being replaced by Solomon Islands Pijin, according to Wikipedia: \url{https://en.wikipedia.org/wiki/Roviana_language}

\section{Digital and Audio Resources}
No audio, video or corpus resources specific to Simbo were found. The Palmer dictionary is the main digital resource.

\section{Curriculum Development Context}
With about 2,700 speakers and little published material, a Simbo curriculum would have to be built almost from scratch. The Palmer dictionary and the old wordlists give a lexical base, but there is no grammar or standard orthography that I could find. Community-led programmes like the Kulu Language Institute are a possible model for early literacy work.
\end{document}